\section{Discussion}\label{sec:discussion}

\paragraph{Guide, do not prune.}
Under capacity coupling, the information a ranking carries is better spent deciding \emph{where
to search} than \emph{what to remove}. Pruning is irreversible and has to keep slack for every
capacity constraint; guidance keeps the full domain reachable and only changes the order in which
the search visits it. This matches the trend from learned pruning towards trust regions
\citep{han2023}, learned neighborhoods \citep{huang2023} and learned delegation
\citep{li2021delegate}. \todo{confirm or revise with the closed-test result}.

\paragraph{The mechanism, not the learning.}
In the pilots the gain over the full solver came from the two restriction mechanisms, and
replacing each learned component by a learning-free one did not reduce it: the LP score in place
of the GNN, rotation in place of the learned selector. We see two reasons. The strong LP already
encodes most of what the GNN can learn about which facilities are worth opening, and it does so
for the instance at hand rather than for a training distribution. And the selector operates on a
set of neighborhoods whose joint local optimum is reached quickly: once every candidate has
failed, no ordering helps. The repetition rate we measured (54--58\% without memory, 25--27\% of
iterations fully exhausted with it) is a direct, inexpensive diagnostic of this situation, and we
suggest reporting it for any LNS before investing in a learned selector.

\paragraph{The classical counterpart is the baseline that matters.}
In every line of work we reconstructed, a classical rule with the same function captured most of
the gain: the LP ranking for pruning, the add-gain filter for swaps, Mettu--Plaxton with local
search for uniform facility location, and a linear continuous approximation for location-routing
(Appendix~\ref{app:reconstructions}). Reporting a learned method only against the raw solver
overstates its contribution.

\paragraph{Low prediction error is not decision quality.}
A routing-cost surrogate with 7.5\% mean absolute percentage error ranked competing
configurations poorly (Kendall $\tau=0.34$), and the optimizer gravitated to configurations where
the surrogate was relatively optimistic.

\paragraph{Timing is the fragile part.}
Two of the errors found by the audit of our own harness affected only timing metrics, yet would
have changed the comparison of anytime behavior. Pinning cores, limiting threads and recording CPU
time next to wall time are cheap and make contention visible.

\paragraph{Limitations.}
\begin{itemize}
\item Most instances are synthetic. Holmberg provides proven optima but small sizes, and the
Olist case relies on assumed capacities and fixed costs.
\item Outside Holmberg, quality is measured against the best known solution.
\item All experiments use one machine with three concurrent pinned runs.
\item The reconstructions in the appendix are methodological reconstructions, not reproductions
of the original code.
\end{itemize}
\todo{add any limitation revealed by the closed test}.
